\section*{РЕЗУЛЬТАТИ ДОСЛІДЖЕННЯ}

Запропоновані механізми апробовано на відтворюваному синтетичному бенчмарку, що моделює узагальнений середній український ЗВО (150~НПП, 20~кафедр; без реального закладу та персональних даних). Якість оцінено за процедурою Монте-Карло ($K=30$ незалежних початкових значень) із наведенням 95\,\% довірчих інтервалів за $t$-розподілом; базові методи й абляції застосовано на \textit{ідентичних} вхідних даних. Повний масив кількісних результатів --- зведені таблиці метрик, аналіз чутливості та оцінку ефективності --- винесено в \textbf{Додаток~А}; нижче узагальнено основні висновки.

\textit{Якість дедуплікації.} Повний конвеєр «золотого запису» статистично відокремлювано домінує над базовими методами: збалансована міра сягає $F_1 = 76{,}7\,\%$ проти $59{,}9\,\%$ у DOI-only та $53{,}5\,\%$ у Title$+$Year (табл.~\ref{tab:res-dedup}, рис.~\ref{fig:res-method}, Додаток~А). Виграш зумовлений насамперед повнотою ($R = 67{,}8\,\%$ проти $42{,}8\,\%$ у DOI-only) за збереження високої точності ($P = 88{,}3\,\%$): за відсутності DOI --- а він відсутній у значній частці записів ORCID та україномовних соціогуманітарних публікацій --- рішення підтримують комплементарні ознаки (близькість назв, рік, перетин множин авторів і видань) у єдиній логістичній моделі $\sigma(\boldsymbol{\beta}\cdot\boldsymbol{\phi})$. Перевага підтверджується й кластерною метрикою $B^3 F_1$ ($95{,}8\,\%$).

\textit{Робастність і передбачувана деградація.} Аналіз чутливості (Додаток~А, рис.~\ref{fig:res-dup}, \ref{fig:res-corr}) засвідчує, що відносне впорядкування методів зберігається на всьому діапазоні частки дублікатів (10--30\,\%), а за зростання спотворення метаданих система деградує \textit{передбачувано}: спершу зростає частка пропущених злиттів (падає повнота), тоді як точність зберігається завдяки охоронцеві конфлікту DOI. Для звітності це бажана властивість --- хибне об'єднання різних праць коштує дорожче за пропуск.

\textit{Якість полів і роль виправленого нормалізатора.} Точність назви злитого запису сягає $99{,}95\,\%$ (Додаток~А, табл.~\ref{tab:res-fields}), що безпосередньо демонструє ефект виправленого критичного дефекту канонізації: попередня реалізація під час зняття пунктуації знищувала \textit{всі кириличні символи}, через що зіставлення за назвою було повністю зламаним для україномовних публікацій. Саме відновлення працездатності кириличного зіставлення є передумовою досягнутих $67{,}8\,\%$ повноти.

\textit{Ефективність та обчислювальна ощадливість.} Рушій виконує повний цикл дедуплікації та злиття для всього модельного ЗВО за $0{,}80 \pm 0{,}08$\,с, тоді як прозора аналітична модель ручної праці дає $\approx 220$\,год/рік на заклад ($\approx 11$\,год на кафедральний звіт, що узгоджується з незалежною оцінкою 8--12\,год); у смугу ручної перевірки (HITL) виноситься $\approx 55$\,пар на дослідника (Додаток~А, табл.~\ref{tab:res-eff}). У термінах рамки TRACE метрикоутворювальні операції виконано детерміновано, без звернення до мовної моделі, що відповідає коефіцієнту обчислювальної ощадливості $\mathrm{CPR}\approx 1$ для ядра консолідації.

\subsection*{Загрози валідності}

Інтерпретацію наведених результатів обмежують кілька загроз валідності, які ми фіксуємо явно.

\textit{Конструктна валідність.} Розподіли спотворень метаданих --- частки втрати DOI, ймовірності перекладу й транслітерації назв, величини зсуву року --- є \textit{модельними припущеннями} про реальну неоднорідність джерел ORCID/OpenAlex/Scopus, а не параметрами, виміряними безпосередньо з них. Відтак абсолютні значення метрик ($F_1 = 76{,}7\,\%$ тощо) слід розглядати як орієнтовні; робастними є \textit{відносні} порівняння методів на ідентичних входах і \textit{тренди} чутливості. Оцінка ручного часу є аналітичною моделлю, а не вимірюванням за участю людей.

\textit{Внутрішня валідність.} Інжектована карта дублікатів слугує gold-розміткою, тож метрики обчислюються відносно \textit{відомої} істини, що виключає помилку анотування. Усі методи й абляції застосовано до спільних входів за фіксованих початкових значень, тому спостережувані відмінності зумовлені саме алгоритмом, а не варіацією даних. Усереднення за 30 початковими значеннями з 95\,\% довірчими інтервалами обмежує вплив випадковості генерації.

\textit{Зовнішня валідність.} Апробацію проведено на \textit{одному} узагальненому модельному ЗВО з фіксованими структурними параметрами; перенесення висновків на установи іншого масштабу чи профілю потребує окремої перевірки. Передусім бракує підтвердження на \textit{реальній} gold-вибірці одного ЗВО з ручною розміткою істинних дублікатів --- цей крок винесено в подальшу роботу як вирішальний для калібрування абсолютних значень. Зазначимо також методологічно істотну обставину: формалізм розв'язання сутностей \textit{реалізовано в коді} рушія (рушій \textit{є} статтею), завдяки чому результати відтворювані безпосередньо запуском конвеєра на синтетичному бенчмарку.
